\documentclass[10pt]{beamer-control}

% Week 10 revision deck: a single standalone presentation summarising and
% linking all nine topics. Built with `make revision`.

\title{\textsc{\large engm\,x304}\\\bigskip {\huge \emph{Applied Control Systems}}\\\bigskip \textsc{Week 10 Revision}}

\usetikzlibrary{fit,backgrounds}

%%%%%%%%%%%%%%%%%%%%%%%%%%%%%%%%
%% REVISION DECK HELPERS
%%%%%%%%%%%%%%%%%%%%%%%%%%%%%%%%

% Header text before the first module
\gdef\currMODULE{Revision}
\gdef\currTOPIC{Course overview}

% "Where did this come from / where does it go" box at the foot of a frame
\newcommand\LINKS[2]{%
  \par\vfill
  {\scriptsize\setlength\fboxsep{3pt}%
  \colorbox{AUHighlight!10}{\parbox{\dimexpr\linewidth-6pt\relax}{%
    \textbf{\color{AUHighlight}$\blacktriangleleft$~Builds on:} #1\par
    \textbf{\color{AUHighlight}$\blacktriangleright$~Leads to:} #2%
    \hfill\hyperlink{coursemap}{\beamerbutton{Course map}}}}}%
}

% Section divider for each module, showing the course map with the module's topics highlighted
\newcommand\MODULEFRAME[2]{%
  \begin{frame}[plain]
    \centering
    {\usebeamerfont{part name}\usebeamercolor[fg]{part name}\partname~\insertromanpartnumber}\par
    \begin{beamercolorbox}[sep=6pt,center]{part title}
      \usebeamerfont{part title}\currMODULE\par
    \end{beamercolorbox}
    \emph{#1}\par\bigskip
    \COURSEMAP{#2}
  \end{frame}
}

% Small recap box used on the topic summary frames
\newcommand\KEYIDEA[1]{%
  {\setlength\fboxsep{4pt}%
  \colorbox{NorthTce!12}{\parbox{\dimexpr\linewidth-8pt\relax}{\small\textbf{Key idea:} #1}}}\par\smallskip
}

\tikzset{
  mapnode/.style = {draw=NorthTce!60, rounded corners=3pt, fill=white,
                    text width=2.7cm, minimum height=1.0cm, align=center,
                    font=\scriptsize, inner sep=2pt},
  maphl/.style   = {fill=NorthTce!25, draw=NorthTce, very thick},
  mapdim/.style  = {},
  maplink/.style = {-{Latex[length=2mm]}, NorthTce!70, thick},
  mapxlink/.style= {-{Latex[length=2mm]}, gray!70, densely dashed},
  t1/.style={mapdim}, t2/.style={mapdim}, t3/.style={mapdim},
  t4/.style={mapdim}, t5/.style={mapdim}, t6/.style={mapdim},
  t7/.style={mapdim}, t8/.style={mapdim}, t9/.style={mapdim},
}

% \COURSEMAP{1,2,3} highlights topics 1-3; \COURSEMAP{} highlights nothing
\newcommand\COURSEMAP[1]{%
  \begingroup
  \foreach \n in {#1} {\global\expandafter\let\csname maphl@\n\endcsname\relax}%
  \def\T##1{\ifcsname maphl@##1\endcsname maphl\else mapdim\fi}%
  \begin{tikzpicture}[node distance=4mm and 9mm]
    \node[mapnode,\T{1}] (T1) {\hyperlink{topic1}{\textbf{1}~Introduction \& System Modelling}};
    \node[mapnode,\T{2},below=of T1] (T2) {\hyperlink{topic2}{\textbf{2}~Dynamic Behaviour}};
    \node[mapnode,\T{3},below=of T2] (T3) {\hyperlink{topic3}{\textbf{3}~Linear Systems}};
    \node[mapnode,\T{4},right=of T1] (T4) {\hyperlink{topic4}{\textbf{4}~State Feedback \& Observers}};
    \node[mapnode,\T{5},below=of T4] (T5) {\hyperlink{topic5}{\textbf{5}~Transfer Functions}};
    \node[mapnode,\T{6},below=of T5] (T6) {\hyperlink{topic6}{\textbf{6}~Frequency Domain Analysis}};
    \node[mapnode,\T{7},right=of T4] (T7) {\hyperlink{topic7}{\textbf{7}~PID Control}};
    \node[mapnode,\T{8},below=of T7] (T8) {\hyperlink{topic8}{\textbf{8}~Frequency Domain Design}};
    \node[mapnode,\T{9},below=of T8] (T9) {\hyperlink{topic9}{\textbf{9}~Robust Performance \& Limits}};
    \node[above=1mm of T1, font=\scriptsize\bfseries, NorthTce] {I. Dynamical Systems};
    \node[above=1mm of T4, font=\scriptsize\bfseries, NorthTce] {II. Control Concepts};
    \node[above=1mm of T7, font=\scriptsize\bfseries, NorthTce] {III. Control Design};
    \draw[maplink] (T1) -- (T2);
    \draw[maplink] (T2) -- (T3);
    \draw[maplink] (T4) -- (T5);
    \draw[maplink] (T5) -- (T6);
    \draw[maplink] (T7) -- (T8);
    \draw[maplink] (T8) -- (T9);
    \draw[maplink] (T3.east) -- (T4.south west);
    \draw[maplink] (T3.east) -- (T5.west);
    \draw[maplink] (T6.east) -- (T8.south west);
    \draw[maplink] (T4.east) -- (T7.west);
    \draw[mapxlink] (T6.south) to[bend right=20] node[below, font=\tiny, text=black]{Nyquist $\to$ robustness} (T9.south);
  \end{tikzpicture}%
  \endgroup
}

% The one block diagram the whole course hangs off (A&M Fig 12.1 notation)
\newcommand\LOOPDIAGRAM{%
  \begin{tikzpicture}[>=Latex, thick, font=\small,
      blk/.style={draw, minimum width=9mm, minimum height=7mm, fill=white},
      sum/.style={draw, circle, inner sep=0pt, minimum size=4mm}]
    \node (r) {$r$};
    \node[blk, right=5mm of r] (F) {$F$};
    \node[sum, right=6mm of F] (s1) {};
    \node[blk, right=6mm of s1] (C) {$C$};
    \node[sum, right=8mm of C] (s2) {};
    \node[blk, right=6mm of s2] (P) {$P$};
    \node[sum, right=8mm of P] (s3) {};
    \node[right=6mm of s3] (y) {$y$};
    \node[above=5mm of s2] (v) {$v$};
    \node[above=5mm of s3] (w) {$w$};
    \draw[->] (r) -- (F);
    \draw[->] (F) -- (s1);
    \draw[->] (s1) -- node[above]{$e$} (C);
    \draw[->] (C) -- node[above]{$u$} (s2);
    \draw[->] (s2) -- (P);
    \draw[->] (P) -- node[above]{$\eta$} (s3);
    \draw[->] (s3) -- (y);
    \draw[->] (v) -- (s2);
    \draw[->] (w) -- (s3);
    \draw[->] ($(s3)!0.5!(y)$) -- ++(0,-1.0) -| node[pos=0.95, left]{$-1$} (s1);
    \node[below=1mm of C, font=\scriptsize, NorthTce] {controller};
    \node[below=1mm of P, font=\scriptsize, NorthTce] {process};
  \end{tikzpicture}%
}

\begin{document}

%%%%%%%%%%%%%%%%%%%%%%%%%%%%%%%%
%% OVERVIEW
%%%%%%%%%%%%%%%%%%%%%%%%%%%%%%%%

\begin{frame}[plain]
\centerline{\includegraphics{au-logo}}
\maketitle
\end{frame}

\begin{frame}{How to use this deck}
\begin{itemize}
\item One pass through the whole course: Modules I $\to$ II $\to$ III
\item Each topic gets a short recap of the ideas and equations you need to know
\item The box at the foot of each recap shows where the idea \emph{came from} and where it is \emph{used next}
\item The course map is clickable: jump to any topic and come back with \beamerbutton{Course map}
\item At the end:
  \begin{itemize}
  \item the same ideas seen in three `languages' (state space, transfer function, frequency response)
  \item running examples traced through the course
  \item a design workflow, a formula sheet, common pitfalls, and self-test questions
  \end{itemize}
\end{itemize}
\bigskip
\alert{This deck is a map, not a replacement.} Use the topic slides, worked examples and \astrom[] for detail.
\end{frame}

\begin{frame}{The question behind the whole course}
\begin{center}
\large\emph{How do we make a dynamical system do what we want,\\ reliably, despite disturbances, noise and model error?}
\end{center}
\bigskip
\begin{columns}[t]
\column{0.32\linewidth}
\textbf{\color{NorthTce}Module I}\\ \emph{What does the system do on its own?}\par\smallskip
\scriptsize Models, states, equilibria, stability, linear response
\column{0.32\linewidth}
\textbf{\color{NorthTce}Module II}\\ \emph{How does feedback change it, and how do we analyse that?}\par\smallskip
\scriptsize State feedback, observers, transfer functions, Bode, Nyquist, margins
\column{0.32\linewidth}
\textbf{\color{NorthTce}Module III}\\ \emph{How do we design a controller that survives reality?}\par\smallskip
\scriptsize PID, sensitivity, loop shaping, root locus, uncertainty, limits
\end{columns}
\end{frame}

\begin{frame}[label=coursemap]{Course map}
\centering
\COURSEMAP{}
\par\medskip
{\scriptsize Linear models (T3) feed both state-space (T4) and transfer-function (T5) methods; pole placement (T4) returns in PID (T7); margins (T6) become sensitivity (T8) and robustness (T9). Click a topic to jump to its recap.\par}
\end{frame}

\begin{frame}{The one diagram to remember}
\centering
\LOOPDIAGRAM
\par\bigskip
\begin{columns}[t]
\column{0.48\linewidth}
\footnotesize
\begin{itemize}
\item $r$ reference, $y$ measured output, $\eta$ true output
\item $v$ load disturbance (usually low frequency)
\item $w$ measurement noise (usually high frequency)
\item $C$ feedback controller, $F$ feedforward/prefilter
\end{itemize}
\column{0.48\linewidth}
\footnotesize
\begin{itemize}
\item Module I: understanding $P$
\item Module II: what $C$ (or $K$, $L$) does to the loop
\item Module III: choosing $C$ and $F$, and how the loop responds to $r$, $v$, $w$ and to errors in $P$
\end{itemize}
\end{columns}
\end{frame}

%%%%%%%%%%%%%%%%%%%%%%%%%%%%%%%%
%% MODULE I
%%%%%%%%%%%%%%%%%%%%%%%%%%%%%%%%

\MODULE{Dynamical Systems}
\MODULEFRAME{What does the system do on its own?}{1,2,3}

\TOPIC{Introduction and System Modelling}

\begin{frame}[label=topic1]{Topic 1 recap: models}
\KEYIDEA{A dynamical system has \emph{states} that evolve in time. Write it as a first-order vector ODE, and control means adding sensors, a controller and actuators to change those dynamics.}
\begin{align*}
\Deriv{x}{t} &= f(x,u) & y &= h(x,u) & &\text{(standard control equation)}\\
\dot x &= Ax + Bu & y &= Cx + Du & &\text{(linear time-invariant)}
\end{align*}
\begin{columns}[t]
\column{0.5\linewidth}
\footnotesize
Mass--spring--damper $m\ddot q + c\dot q + kq = f$:
\begin{align*}
x &= \Matr{q\\\dot q}, &
A &= \Matr{0 & 1\\ -k/m & -c/m}, &
B &= \Matr{0\\1/m}
\end{align*}
\column{0.5\linewidth}
\footnotesize
\begin{itemize}
\item Feedback: act on measured output; feedforward: act on known inputs
\item Why feedback: robustness to uncertainty, design of dynamics
\item Costs: stability, coupling, complexity
\item Discrete time: $x[k{+}1] = Ax[k] + Bu[k]$
\end{itemize}
\end{columns}
\LINKS{(start of the course)}{Every later topic uses $\dot x = f(x,u)$ or $\dot x = Ax+Bu$. Difference equations return in PID implementation (T7).}
\end{frame}

\begin{frame}{Topic 1 recap: models from experiments}
\begin{columns}[t]
\column{0.5\linewidth}
\footnotesize
\textbf{Step-response models}
\begin{itemize}
\item ITD: integrator + time delay ($a$, $\TimeDel$)
\item FOTD: $\StaticGain$, $\TimeDel$, $\TimeConst$ (63\% rise)
\item SOTD: $\StaticGain$, overshoot, peak time, rise/settling time
\end{itemize}
\textbf{Normalise and scale} inputs/outputs by $r_{\max}$, $y_{\max}$ so results generalise.
\column{0.5\linewidth}
\footnotesize
\textbf{Second order} ($m\ddot q + c\dot q + kq = f$)
\begin{align*}
\wn &= \sqrt{k/m} & \zz &= \frac{c}{2\sqrt{km}}\\
\wdamp &= \wn\sqrt{1-\zz^2} & \wdamp &\approx \frac{2\pi}{T_2-T_1}
\end{align*}
Decay ratio of successive peaks $d$:
\begin{align*}
\zz \approx \frac{1}{\sqrt{1+\left(2\pi/\ln d\right)^2}}
\end{align*}
\end{columns}
\LINKS{Step response features (T1).}{ITD/FOTD models are exactly what Ziegler--Nichols, CHR and Cohen--Coon tuning use (T7). $\zz$, $\wn$ return in pole placement (T4).}
\end{frame}

\TOPIC{Dynamic Behaviour}

\begin{frame}[label=topic2]{Topic 2 recap: qualitative behaviour and stability}
\KEYIDEA{Close the loop with $u=\alpha(x)$ and you get $\dot x = F(x)$. Its equilibria and their stability tell you most of what matters.}
\begin{columns}[t]
\column{0.5\linewidth}
\footnotesize
\begin{itemize}
\item Equilibrium: $F(x_e) = 0$
\item Vector field / phase portrait: pictures of $F(x)$ and of trajectories
\item Limit cycle: isolated closed trajectory (oscillators, friction)
\item Nonlinear surprises: finite escape time ($\dot x = x^2$), multiple equilibria (pendulum up/down)
\end{itemize}
\column{0.5\linewidth}
\footnotesize
\begin{itemize}
\item \textbf{Lyapunov stable}: start close, stay close
\item \textbf{Asymptotically stable}: and converge
\item \textbf{Local} vs \textbf{global} (ball $B_r(a)$)
\item Lyapunov function $V(x)$: energy-like; $\dot V<0 \Rightarrow$ stable
\item \alert{Linear: stable $\iff \Real\lambda_j(A)<0$ for all $j$}, and stability is global
\end{itemize}
\end{columns}
\smallskip
\footnotesize Running examples introduced here: \textbf{cruise control} (nonlinear model + PI) and \textbf{inverted pendulum / balance system}.
\LINKS{State models (T1).}{Eigenvalue test justified via the matrix exponential (T3). Stability of the \emph{closed loop} judged by Nyquist (T6).}
\end{frame}

\TOPIC{Linear Systems}

\begin{frame}[label=topic3]{Topic 3 recap: linear systems}
\KEYIDEA{Linearity gives superposition, so the whole response follows from the matrix exponential, and $\ee^{At}$ is built from the modes $\ee^{\lambda_j t}$.}
\vspace{-1ex}
\begin{align*}
x(t) &= \ee^{At}x(0) + \int_0^t \ee^{A(t-\tau)}Bu(\tau)\,\dee\tau & &\text{(convolution equation)}\\
\ee^{At} &= I + At + \tfrac{1}{2}A^2t^2 + \cdots & \ee^{\lambda t} &= \ee^{\sigma t}(\cos\ww t + \ii\sin\ww t)
\end{align*}
\begin{columns}[t]
\column{0.5\linewidth}
\footnotesize
\begin{itemize}
\item LTI: superposition + time invariance
\item Diagonal/Jordan form $J = T^{-1}AT$ exposes the eigenvalues
\item Coordinate changes do not change eigenvalues
\end{itemize}
\column{0.5\linewidth}
\footnotesize
\begin{itemize}
\item Step steady state: $D - CA^{-1}B$
\item Frequency response: $G(\ii\ww) = C(\ii\ww I-A)^{-1}B + D$
\item Gain $|G(\ii\ww)|$, phase $\arg G(\ii\ww)$, bandwidth $\wb$, peak $\Mpeak$
\end{itemize}
\end{columns}
\LINKS{Stability via eigenvalues (T2).}{$G(s)=C(sI-A)^{-1}B+D$ becomes the transfer function (T5); $D-CA^{-1}B$ becomes DC gain $G(0)$ (T5).}
\end{frame}

\begin{frame}{Topic 3 recap: linearisation}
\begin{columns}[t]
\column{0.55\linewidth}
\footnotesize
\textbf{Jacobian linearisation} about $(x_e,u_e)$:
\begin{align*}
z &= x-x_e, \quad v = u-u_e, \quad w = y - h(x_e,u_e)\\
\dot z &= Az + Bv, \qquad w = Cz + Dv\\
A &= \left.\PDeriv{f}{x}\right|_{e} \quad B = \left.\PDeriv{f}{u}\right|_{e} \quad
C = \left.\PDeriv{h}{x}\right|_{e} \quad D = \left.\PDeriv{h}{u}\right|_{e}
\end{align*}
\begin{itemize}
\item Always shift to the equilibrium first (linear, not affine)
\item Valid only near the operating point
\end{itemize}
\column{0.42\linewidth}
\footnotesize
\textbf{Feedback linearisation}
\begin{itemize}
\item Cancel a \emph{known} nonlinearity: $\dot x = x^2+u$, $u = -x^2-kx$
\item Dynamic inversion: $u = B^{-1}(Mv + C)$ gives $\ddot q = v$
\item Relies on an accurate model and capable actuators
\end{itemize}
\end{columns}
\bigskip
\footnotesize
Example: pendulum $\dot x = \Matr{x_2\\ \sin x_1 - cx_2}$. Upright: $A=\Matr{0&1\\1&-c}$, one eigenvalue $>0$: \alert{unstable}. Hanging: $A=\Matr{0&1\\-1&-c}$: stable for $c>0$.
\LINKS{Nonlinear models (T1, T2).}{Every linear tool in Modules II and III assumes this step has been done. Parameters depending on the operating point become \emph{parametric uncertainty} (T9).}
\end{frame}

\begin{frame}{Bridge I $\to$ II: everything is in the eigenvalues}
\begin{theorem}
The system $\dot x = Ax$ is asymptotically stable if and only if all eigenvalues of $A$ have strictly negative real part.
\end{theorem}
\medskip
Write each eigenvalue as $\lambda = \sigma \pm \ii\ww$:
\begin{itemize}
\item $\sigma$ sets growth/decay $\to$ \emph{stability} and \emph{settling time}
\item $\ww$ sets oscillation $\to$ \emph{damped frequency}
\item the angle from the negative real axis sets \emph{damping ratio}: $\zz = \cos\varphi$
\end{itemize}
\bigskip
So the natural next questions are:
\begin{enumerate}
\item Can feedback \emph{move} the eigenvalues? \hfill \alert{$\to$ T4 state feedback}
\item Is there a simpler description of the input/output map? \hfill \alert{$\to$ T5 transfer functions}
\item Can we judge closed-loop stability without computing eigenvalues? \hfill \alert{$\to$ T6 Nyquist}
\end{enumerate}
\end{frame}

%%%%%%%%%%%%%%%%%%%%%%%%%%%%%%%%
%% MODULE II
%%%%%%%%%%%%%%%%%%%%%%%%%%%%%%%%

\MODULE{Control System Concepts}
\MODULEFRAME{How does feedback change the system, and how do we analyse it?}{4,5,6}

\TOPIC{State Feedback}

\begin{frame}[label=topic4]{Topic 4 recap: reachability and state feedback}
\KEYIDEA{If the system is reachable, state feedback $u=-Kx+k_f r$ can put the closed-loop eigenvalues anywhere.}
\begin{columns}[t]
\column{0.5\linewidth}
\footnotesize
\textbf{Reachability}
\begin{align*}
W_r = \Matr{B & AB & \cdots & A^{n-1}B}
\end{align*}
Reachable $\iff W_r$ full rank (invertible for single input). For linear systems, reachable $\iff$ controllable.
\par\medskip
\textbf{Closed loop}
\begin{align*}
\dot x = (A-BK)x + Bk_f r
\end{align*}
\column{0.5\linewidth}
\footnotesize
\textbf{Eigenvalue assignment}
\begin{enumerate}
\item $\det(sI-A) = s^n + a_1 s^{n-1} + \cdots + a_n$
\item Choose $p(s) = s^n + p_1 s^{n-1} + \cdots + p_n$
\item $K = \Matr{p_1-a_1 & \cdots & p_n-a_n}\tilde W_r W_r^{-1}$
\item $k_f$ for unity DC gain:\\ $k_f = -1/\bigl(C(A-BK)^{-1}B\bigr)$
\end{enumerate}
Matlab: \texttt{place}, \texttt{acker}
\end{columns}
\LINKS{$A$, $B$, eigenvalues and $D-CA^{-1}B$ (T3).}{Same pole-placement idea for PI/PID (T7). Placing poles badly hurts robustness (T9).}
\end{frame}

\begin{frame}{Topic 4 recap: choosing poles, integral action}
\begin{columns}[t]
\column{0.5\linewidth}
\footnotesize
\textbf{Second-order design targets}\\
$\ddot q + 2\zz\ww_0\dot q + \ww_0^2 q = k\ww_0^2 u$,\quad $\lambda = -\zz\ww_0 \pm \ww_0\sqrt{\zz^2-1}$
\par\smallskip
\renewcommand\arraystretch{1.3}
\begin{tabular}{@{}ll@{}}
\toprule
Overshoot & $M_p = \ee^{-\pi\zz/\sqrt{1-\zz^2}}$\\
2\% settling & $T_s \approx 4/(\zz\ww_0)$\\
Rise time & $T_r = \ee^{\varphi/\tan\varphi}/\ww_0$, $\varphi=\arccos\zz$\\
\bottomrule
\end{tabular}
\par\smallskip
Higher order: place the \alert{dominant} (slowest) pair for the specs, the rest faster.
\column{0.5\linewidth}
\footnotesize
\textbf{Integral action}: augment with $z = \int (y-r)\,\dee t$
\begin{align*}
\Deriv{}{t}\Matr{x\\z} &= \Matr{Ax+Bu\\ Cx-r}\\
u &= -Kx - \kki z + k_f r
\end{align*}
At steady state $\dot z = 0 \Rightarrow y = r$, even with model error or constant disturbances.
\par\smallskip
Cruise control, $P = b/(s+a)$:
\begin{align*}
\kkp = \frac{a_1 - a}{b}, \quad \kki = \frac{a_2}{b}
\end{align*}
\end{columns}
\LINKS{Second-order response (T1, T3).}{The \emph{same} $\kkp$, $\kki$ formulas reappear as PI pole placement (T7). Integral action $\leftrightarrow$ $1/s$ in $C$ $\leftrightarrow$ $G_{yv}\approx s/\kki$ (T8).}
\end{frame}

\begin{frame}{Topic 4 recap: observers and output feedback}
\begin{columns}[t]
\column{0.5\linewidth}
\footnotesize
\textbf{Observability}
\begin{align*}
\Wob = \Matr{C\\CA\\\vdots\\CA^{n-1}}
\end{align*}
Observable $\iff \Wob$ has rank $n$.
\par\medskip
\textbf{Observer}
\begin{align*}
\Deriv{\hat x}{t} &= A\hat x + Bu + L(y - C\hat x)\\
\Deriv{\tilde x}{t} &= (A-LC)\tilde x, \quad \tilde x = x - \hat x
\end{align*}
\column{0.5\linewidth}
\footnotesize
\textbf{Duality}: $A\leftrightarrow A\Tr$, $B\leftrightarrow C\Tr$, $K\leftrightarrow L\Tr$, $W_r\leftrightarrow \Wob\Tr$\\
(same \texttt{place}/\texttt{acker} algorithms)
\par\medskip
\textbf{Controller + observer}, $u = -K\hat x + k_f r$:
\begin{align*}
\Deriv{}{t}\Matr{x\\\tilde x} = \Matr{A-BK & BK\\ 0 & A-LC}\Matr{x\\\tilde x} + \Matr{Bk_f\\0}r
\end{align*}
Block triangular, so
\begin{align*}
p(s) = \det(sI-A+BK)\,\det(sI-A+LC)
\end{align*}
\end{columns}
\LINKS{Reachability and eigenvalue assignment (this topic).}{An observer-based controller is just a transfer function $C(s)$ (T5); its robustness is checked with the Gang of Four (T8, T9).}
\end{frame}

\TOPIC{Transfer Functions}

\begin{frame}[label=topic5]{Topic 5 recap: transfer functions}
\KEYIDEA{For input $\ee^{st}$, an LTI system returns $G(s)\ee^{st}$. The transfer function $G(s)$ is unique even though $A,B,C,D$ are not.}
\vspace{-1ex}
\begin{align*}
G(s) = C(sI-A)^{-1}B + D = \frac{b(s)}{a(s)} = \frac{b_0 s^m + \cdots + b_m}{s^n + a_1 s^{n-1} + \cdots + a_n}
\end{align*}
\begin{columns}[t]
\column{0.5\linewidth}
\footnotesize
\begin{itemize}
\item \alert{Poles} = roots of $a(s)$ = eigenvalues of $A$ (modes $\ee^{pt}$)
\item \alert{Zeros} = roots of $b(s)$ (blocked signals $\ee^{qt}$)
\item Pole excess $\npoleexcess = n-m$; proper if $\npoleexcess\ge 0$
\item Minimal realisation removes cancelled pole/zero pairs
\end{itemize}
\column{0.5\linewidth}
\footnotesize
\textbf{Laplace essentials}
\begin{align*}
\Lapl\{\dot f\} &= sF(s) & \Lapl\left\{\textstyle\int f\right\} &= F(s)/s\\
f(\infty) &= \lim_{s\to 0} sF(s) & f(0^+) &= \lim_{s\to\infty} sF(s)
\end{align*}
Step input $R=1/s$: \alert{$y(\infty) = G(0)$}, the DC gain
\begin{align*}
G(0) = \frac{b_m}{a_n} = D - CA^{-1}B
\end{align*}
\end{columns}
\LINKS{Frequency response and steady state (T3).}{Loop transfer function $L=PC$ (T6); Gang of Four (T8).}
\end{frame}

\begin{frame}{Topic 5 recap: block diagrams and pole/zero cancellation}
\begin{columns}[t]
\column{0.5\linewidth}
\footnotesize
\textbf{Feedback loop} (negative feedback, $F=1$):
\begin{align*}
G_{yr} = \frac{PC}{1+PC}, \quad
G_{yv} = \frac{P}{1+PC}, \quad
G_{yw} = \frac{1}{1+PC}
\end{align*}
Block algebra: series multiply, parallel add, loop $\frac{\text{forward}}{1+\text{loop}}$.
\par\medskip
\textbf{Algebraic loops}: a direct feedthrough ($D\neq 0$) inside a loop needs an iterative solver in simulation.
\column{0.5\linewidth}
\footnotesize
\textbf{Pole/zero cancellation trap}
\begin{align*}
P &= \frac{s-1}{(s+1)(s+5)}, & C &= \frac{s+1}{(s-1)(s+2)}
\end{align*}
$PC$ looks stable, but:
\begin{itemize}
\item the controller itself is unstable
\item cancellation is never exact in practice
\item the unstable pole is still visible in other paths, e.g.\ $G_{yv}$
\end{itemize}
\alert{Never cancel RHP poles or zeros.}
\end{columns}
\LINKS{TF algebra (this topic).}{All six closed-loop TFs (Gang of Six) and internal stability (T8).}
\end{frame}

\begin{frame}{Topic 5 recap: Bode plots and time delays}
\begin{columns}[t]
\column{0.5\linewidth}
\footnotesize
\textbf{Sketching rules}
\begin{itemize}
\item Start at DC gain, slope \SI{0}{dB/dec}
\item Each pole: $-\SI{20}{dB/dec}$, phase $-\ang{90}$ over about a decade
\item Each zero: $+\SI{20}{dB/dec}$, phase $+\ang{90}$
\item Integrator = pole at $0$; differentiator = zero at $0$
\item $\theta_0 = -\ang{90}(N_I-N_D)$, $\theta_\infty = -\ang{90}(N_P-N_Z)$
\end{itemize}
\textbf{dB}: $20\log_{10}|G|$;\; $\times 10 = +\SI{20}{dB}$;\; $\times 2 \approx +\SI{6}{dB}$;\; $1/\sqrt2 \approx -\SI{3}{dB}$
\column{0.5\linewidth}
\footnotesize
\textbf{Time delay} $\ee^{-\tau s}$
\begin{itemize}
\item $|\cdot| = 1$, phase $-\ww\tau$ (unbounded lag)
\item Padé: $\ee^{-\tau s}\approx \dfrac{1-\tau s/2}{1+\tau s/2}$, a \alert{RHP zero} at $2/\tau$
\end{itemize}
\textbf{Experimental FRF}: stepped sine, chirp or noise; \texttt{tfestimate}, \texttt{tfest}
\par\medskip
\textbf{RHP zeros and delays} add phase lag without changing gain: \emph{non-minimum phase}.
\end{columns}
\LINKS{$G(\ii\ww)$ (T3).}{Bode plot of $L$ gives gain and phase margins (T6); loop shaping is done on the Bode plot (T8); delays and RHP zeros limit performance (T9).}
\end{frame}

\TOPIC{Frequency Domain Analysis}

\begin{frame}[label=topic6]{Topic 6 recap: loop transfer function and Nyquist}
\KEYIDEA{Judge \emph{closed-loop} stability from the \emph{open-loop} $L = PC$: a sinusoid that returns around the loop with gain 1 and phase $-\ang{180}$ ($L(\ii\ww)=-1$) is self-sustaining.}
\begin{columns}[t]
\column{0.5\linewidth}
\footnotesize
\textbf{Nyquist criterion}
\begin{align*}
n_Z = n_P + n_N
\end{align*}
\begin{itemize}
\item $n_Z$: RHP closed-loop poles (zeros of $1+L$)
\item $n_P$: RHP open-loop poles of $L$ (known)
\item $n_N$: clockwise encirclements of $-1$
\end{itemize}
Stable $L$ ($n_P=0$): \alert{no encirclements of $-1$}.\\
Unstable $L$: need $n_P$ \emph{counter}-clockwise encirclements.
\column{0.5\linewidth}
\footnotesize
\textbf{Example} $L = \dfrac{k}{s(s+1)^2}$
\begin{itemize}
\item phase $-\ang{180}$ at $\ww=1$, where $L(\ii) = -k/2$
\item encircles $-1$ for $k>2$: unstable
\end{itemize}
\textbf{Delay}: $L\,\ee^{-\tau s}$ keeps $|L|$ but rotates the curve; the critical delay is
\begin{align*}
\tau_c = \frac{\pi + \theta_0}{\ww_0}
\end{align*}
\end{columns}
\LINKS{Eigenvalue stability (T2, T3), $G(\ii\ww)$ and Bode (T5).}{Distance to $-1$ becomes $\Msens$ and sensitivity (T8) and robust stability (T9). ZN frequency tuning uses the $-\ang{180}$ point (T7).}
\end{frame}

\begin{frame}{Topic 6 recap: stability margins}
\begin{columns}[c]
\column{0.5\linewidth}
\includegraphics[width=\linewidth]{figure10.11}
\par{\tiny \AMref*{Fig.\ 10.11}\par}
\column{0.5\linewidth}
\footnotesize
\begin{itemize}
\item \textbf{Gain margin} $\gainmargin$: factor on $|L|$ to reach $-1$ at phase crossover. Aim $2<\gainmargin<5$
\item \textbf{Phase margin} $\phasemargin$: extra lag to reach $-1$ at gain crossover $\ww\gaincrossover$. Aim $\ang{30}<\phasemargin<\ang{60}$
\item \textbf{Stability margin} $\stabilitymargin$: shortest distance from $L(\ii\ww)$ to $-1$
\item \textbf{Delay margin}: $\phasemargin = \ww\gaincrossover\delaymargin$ ($\phasemargin$ in rad)
\end{itemize}
\end{columns}
\medskip
\footnotesize
Example $L = 3/(s+1)^3$: $\stabilitymargin = 0.464$, $\gainmargin = 2.67$, $\phasemargin = \ang{41.7}$.
\par
\textbf{Bode's relation} (minimum phase): slope $n$ on the gain curve $\Rightarrow$ phase $\approx n\cdot\ang{90}$. RHP zeros, RHP poles and delays add extra lag.
\LINKS{Nyquist criterion (this topic), Bode plots (T5).}{$\Msens = 1/\stabilitymargin$ (T8); target slope $-1$ at crossover in loop shaping (T8); margins quantify robustness (T9).}
\end{frame}

\begin{frame}{Bridge II $\to$ III: one fact, three views}
\begin{center}
\small
\renewcommand\arraystretch{1.4}
\setlength\tabcolsep{4pt}\begin{tabular}{@{}>{\raggedright}p{2.3cm}>{\raggedright}p{2.6cm}>{\raggedright}p{2.6cm}>{\raggedright\arraybackslash}p{2.6cm}@{}}
\toprule
& \textbf{State space (T4)} & \textbf{Transfer fn (T5)} & \textbf{Frequency (T6)}\\
\midrule
Closed-loop char.\ polynomial & $\det(sI-A+BK)$ & $a(s)+k\,b(s)$ \newline (num.\ of $1+L$) & zeros of $1+L(s)$\\
Stable when & $\Real\lambda(A-BK)<0$ & CL poles in LHP & $n_Z = n_P + n_N = 0$\\
Steady state & $D - CA^{-1}B$ & $G(0)$, FVT & low-$\ww$ gain\\
How much margin & pole locations & pole locations & $\gainmargin$, $\phasemargin$, $\stabilitymargin$\\
\bottomrule
\end{tabular}
\end{center}
\bigskip
Module III asks: \emph{given these tools, how do we pick a good controller?}
\begin{itemize}
\item a simple, model-free structure \hfill \alert{$\to$ T7 PID}
\item shaping $L$ so the loop handles $r$, $v$, $w$ well \hfill \alert{$\to$ T8 sensitivity, loop shaping}
\item checking it still works when $P$ is wrong \hfill \alert{$\to$ T9 robustness, limits}
\end{itemize}
\end{frame}

%%%%%%%%%%%%%%%%%%%%%%%%%%%%%%%%
%% MODULE III
%%%%%%%%%%%%%%%%%%%%%%%%%%%%%%%%

\MODULE{Control System Design}
\MODULEFRAME{How do we design a controller that survives reality?}{7,8,9}

\TOPIC{PID Control}

\begin{frame}[label=topic7]{Topic 7 recap: PID structure}
\KEYIDEA{P acts on the present error, I on the past, D on the predicted future. Three knobs control most industrial loops.}
\vspace{-1ex}
\begin{align*}
u &= \kkp e + \kki\int_0^t e\,\dee\tau + \kkd\Deriv{e}{t}
= \kkp\Bigl(e + \frac{1}{\TTi}\int_0^t e\,\dee\tau + \TTd\Deriv{e}{t}\Bigr)\\
\TTi &= \kkp/\kki, \qquad \TTd = \kkd/\kkp
\end{align*}
\begin{center}
\footnotesize
\renewcommand\arraystretch{1.2}
\begin{tabular}{@{}lll@{}}
\toprule
Term & Increasing the gain\ldots & Watch out for\ldots\\
\midrule
P & faster, smaller (but nonzero) steady-state error & oscillation, instability\\
I & removes steady-state error, rejects load disturbances & overshoot, windup\\
D & adds damping, reduces overshoot & noise amplification\\
\bottomrule
\end{tabular}
\end{center}
\footnotesize
\begin{itemize}
\item $\kkd s$ is improper: filter it, $\dfrac{\kkd s}{1+s\TTf}$ with $\TTf = \TTd/N$, $5\le N\le 20$
\item PI on $P=b/(s+a)$: $s^2 + (a+b\kkp)s + b\kki$, so $\kkp = (a_1-a)/b$, $\kki = a_2/b$ \alert{(identical to T4)}
\end{itemize}
\LINKS{Integral action and pole placement (T4), simple models (T1).}{PD with filter = lead compensator, PI = lag compensator (T8).}
\end{frame}

\begin{frame}{Topic 7 recap: tuning}
\begin{columns}[t]
\column{0.5\linewidth}
\footnotesize
\textbf{Ziegler--Nichols (1942)}
\begin{itemize}
\item Step method: ITD model $P = \frac{a}{\tau s}\ee^{-\tau s}$
\item Frequency method: raise $\kkp$ until oscillation: critical gain $k_c$, period $T_c$; this is the \alert{$-\ang{180}$ point on the Nyquist curve}
\item Good starting point; often too much overshoot for mechanical systems
\end{itemize}
\textbf{Improved rules on FOTD} $\frac{K}{1+sT}\ee^{-\tau s}$
\begin{itemize}
\item $\tau_n = \tau/(\tau+T)$: lag ($\to 0$) vs delay ($\to 1$) dominated
\item CHR (0\% or 20\% overshoot; set point or disturbance), Cohen--Coon
\end{itemize}
\column{0.5\linewidth}
\scriptsize
\textbf{Ziegler--Nichols rules} \AMref{Table 11.1}\par\smallskip
\begin{tabular}{@{}llll@{}}
\toprule
Step & $\kkp$ & $\TTi$ & $\TTd$\\
\midrule
P & $1/a$ & & \\
PI & $0.9/a$ & $\tau/0.3$ & \\
PID & $1.2/a$ & $\tau/0.5$ & $0.5\tau$\\
\bottomrule
\end{tabular}\par\smallskip
\begin{tabular}{@{}llll@{}}
\toprule
Freq. & $\kkp$ & $\TTi$ & $\TTd$\\
\midrule
P & $0.5k_c$ & & \\
PI & $0.45k_c$ & $T_c/1.2$ & \\
PID & $0.6k_c$ & $T_c/2$ & $T_c/8$\\
\bottomrule
\end{tabular}
\par\medskip
\footnotesize
\textbf{Relay autotuning}: relay amplitude $d$, output amplitude $a$
\begin{align*}
k_c = \frac{4d}{\pi a}, \qquad \ww_c = \frac{2\pi}{T_c}
\end{align*}
\end{columns}
\LINKS{ITD/FOTD step-response models (T1), Nyquist critical point (T6).}{Check the tuned loop's margins and sensitivity peaks (T6, T8).}
\end{frame}

\begin{frame}{Topic 7 recap: windup and implementation}
\begin{columns}[t]
\column{0.5\linewidth}
\footnotesize
\textbf{Integrator windup}
\begin{itemize}
\item Saturated actuator $\Rightarrow$ loop is effectively open
\item Integrator keeps growing $\Rightarrow$ large overshoot, slow recovery
\item Fix: \alert{anti-windup} feedback of $e_s = u_{\text{sat}} - u$ into the integrator with gain $k_{\text{aw}}$ (tracking)
\item Applies to any controller with integral action, including state-space designs
\end{itemize}
\textbf{Setpoint weighting} (2-DOF PID):
\begin{align*}
u = \kkp(\beta r - y) + \kki\!\int\!(r-y) + \kkd\Bigl(\gamma\dot r - \dot y\Bigr)
\end{align*}
Changes response to $r$ only, not to $v$ or $w$.
\column{0.5\linewidth}
\footnotesize
\textbf{Digital implementation} (sample time $h$)
\begin{align*}
P(t_k) &= \kkp\bigl(\beta r(t_k) - y(t_k)\bigr)\\
I(t_{k+1}) &= I(t_k) + \kki h\, e(t_k) + \tfrac{h}{T_{\text{aw}}}\bigl(\mathrm{sat}(u)-u\bigr)\\
D(t_k) &= \tfrac{\TTf}{\TTf+h}D(t_{k-1}) - \tfrac{\kkd}{\TTf+h}\bigl(y(t_k)-y(t_{k-1})\bigr)
\end{align*}
Same idea as the discrete PI in T1:
\begin{align*}
u_i[k+1] = u_i[k] + h\,e[k]
\end{align*}
\end{columns}
\LINKS{Difference equations (T1), saturation in block diagrams (T1).}{Actuator limits as a fundamental performance limit (T9).}
\end{frame}

\TOPIC{Frequency Domain Design}

\begin{frame}[label=topic8]{Topic 8 recap: Gang of Four}
\KEYIDEA{Four transfer functions describe how the feedback loop responds to disturbances and noise. $S+T=1$ means you cannot make both small at the same frequency.}
\begin{columns}[c]
\column{0.5\linewidth}
\footnotesize
\renewcommand\arraystretch{1.5}
\begin{tabular}{@{}ll@{}}
\toprule
$S = \dfrac{1}{1+PC}$ & sensitivity: $v$-rejection, tracking\\
$T = \dfrac{PC}{1+PC}$ & complementary: noise $\to y$\\
$PS = \dfrac{P}{1+PC}$ & load disturbance $v \to y$\\
$CS = \dfrac{C}{1+PC}$ & noise $w \to u$\\
\bottomrule
\end{tabular}
\column{0.5\linewidth}
\footnotesize
\begin{itemize}
\item \textbf{Internally stable} $\iff$ all four are stable (catches hidden cancellations, T5)
\item Gang of Six adds $G_{yr} = TF$ and $G_{ur} = CSF$, which are shaped by the feedforward $F$
\item $|S(\ii\ww)|<1$: disturbance attenuated; $>1$: amplified
\item $\Msens = \max|S| = 1/\stabilitymargin$ \alert{(the Nyquist distance from T6)}
\end{itemize}
\end{columns}
\medskip
\footnotesize
Low frequency: integral action gives $G_{yv}\approx s/\kki \to 0$. High frequency: $-G_{uw}\approx C$, so roll off $C$ (filter D).
\LINKS{Block diagram TFs and cancellation (T5), stability margin (T6).}{Sensitivity to model error $\frac{\dee G}{G} = S\frac{\dee P}{P}$, robust stability $|\delta|<1/|T|$ (T9).}
\end{frame}

\begin{frame}{Topic 8 recap: specifications and loop shaping}
\begin{columns}[t]
\column{0.48\linewidth}
\includegraphics[width=\linewidth]{figure12.8}
\par{\tiny \AMref*{Fig.\ 12.8}\par}
\footnotesize
\begin{itemize}
\item High $|L|$ at low $\ww$: tracking, disturbance rejection
\item Low $|L|$ at high $\ww$: noise immunity, robustness
\item Crossover slope $n\gaincrossover \approx -2 + \phasemargin/\ang{90}$
\end{itemize}
\column{0.5\linewidth}
\footnotesize
\textbf{Specs}: overshoot, rise time, settling time $\leftrightarrow$ peak $M_r$, bandwidth $\wb$, $\ww\gaincrossover$, $\Msens$, $M_t$
\begin{itemize}
\item Rise time--bandwidth: $T_r\wb \approx 3$
\item Large overshoot $\leftrightarrow$ large resonant peak
\item Ramp error $= v_0 e_1$, from $G_{er}\approx e_1 s + e_2 s^2 + \cdots$
\end{itemize}
\textbf{Lead/lag}: $C = k\dfrac{s+a}{s+b}$
\begin{itemize}
\item Lead ($a<b$): adds phase near crossover; = filtered PD
\item Lag ($a>b$): raises low-$\ww$ gain; PI is lag with $b=0$
\end{itemize}
\end{columns}
\LINKS{Bode sketching (T5), Bode's gain--phase relation (T6), PID (T7).}{Bode's integral limits how far $|S|$ can be pushed down (T9).}
\end{frame}

\begin{frame}{Topic 8 recap: root locus}
\begin{columns}[t]
\column{0.5\linewidth}
\footnotesize
Closed-loop poles as one gain $k$ varies, $C = k$:
\begin{align*}
a\closedloop(s) = a(s) + k\,b(s) = 0
\end{align*}
\begin{itemize}
\item $n$ branches start at the open-loop poles ($k=0$)
\item They end at open-loop zeros or go to infinity
\item $\npoleexcess$ asymptotes from the centroid, angles $(\pi + 2l\pi)/\npoleexcess$
\item RHP zeros, or $\npoleexcess \ge 3$: \alert{unstable for large gain}
\end{itemize}
\column{0.5\linewidth}
\footnotesize
\textbf{Example}: $P = 1/(s^2+1)$
\begin{itemize}
\item PI, $C = k\frac{s+1}{s}$: unstable for all $k$
\item PID, $C = k\frac{s^2+2s+2}{s}$: stable, choose $k$ to place poles
\end{itemize}
\textbf{Example}: $L = \dfrac{k(s-2)}{(s+4)^2}$
\begin{itemize}
\item $s^2 + (8+k)s + (16-2k)$: unstable for $k>8$
\item One branch heads to the RHP zero at $s=2$
\end{itemize}
\end{columns}
\LINKS{Characteristic polynomial (T4, T5), PID pole placement (T7).}{RHP zeros as a fundamental limit (T9).}
\end{frame}

\TOPIC{Robust Performance \& Fundamental Limits}

\begin{frame}[label=topic9]{Topic 9 recap: uncertainty and robustness}
\KEYIDEA{Every model is wrong. Feedback tolerates model error where $|S|$ and $|T|$ allow it, and some limits cannot be removed by any controller.}
\begin{columns}[t]
\column{0.5\linewidth}
\footnotesize
\textbf{Uncertainty types}
\begin{itemize}
\item Parametric (e.g.\ vehicle mass in cruise control)
\item Unmodelled dynamics: additive $\Delta$, multiplicative $\delta = \Delta/P$
\item Open-loop similarity can mislead: compare systems in closed loop
\end{itemize}
\textbf{Robust stability}
\begin{align*}
|\delta(\ii\ww)| < \frac{1}{|T(\ii\ww)|}\;\forall\ww
\quad\Leftarrow\quad |\delta| < \frac{1}{M_t}
\end{align*}
\column{0.5\linewidth}
\footnotesize
\textbf{Robust performance}
\begin{align*}
\frac{\dee G_{yr}}{G_{yr}} = S\frac{\dee P}{P}, \qquad
\frac{\dee G_{uw}}{G_{uw}} = -T\frac{\dee P}{P}
\end{align*}
\begin{itemize}
\item Where $|S|$ is small (low $\ww$), the response is insensitive to $P$
\item Where $|T|$ is small (high $\ww$), large relative errors are tolerated
\end{itemize}
\textbf{Robust pole placement}
\begin{itemize}
\item Cancel \emph{slow stable} zeros with controller poles
\item Put controller zeros near \emph{fast stable} poles
\end{itemize}
\end{columns}
\LINKS{Gang of Four (T8), Nyquist distance to $-1$ (T6), pole placement (T4, T7).}{(end of the course)}
\end{frame}

\begin{frame}{Topic 9 recap: fundamental limits}
\begin{columns}[t]
\column{0.5\linewidth}
\footnotesize
\textbf{Stabilisability}
\begin{itemize}
\item Reachable $\Rightarrow$ stabilisable; otherwise the unreachable part must be stable
\item PBH: $\operatorname{rank}\Matr{A-\lambda I & B} = n$ for all $\lambda$ (reachable), for $\Real\lambda\ge 0$ (stabilisable)
\item Strongly stabilisable (stable $C$) $\iff$ an even number of real poles between every pair of real RHP zeros
\end{itemize}
\textbf{RHP zeros and delays}
\begin{itemize}
\item Zeros depend on sensor/actuator placement: sometimes redesign removes them (vehicle steering)
\item Delays act like RHP zeros (Padé zero at $2/\tau$)
\end{itemize}
\column{0.5\linewidth}
\footnotesize
\textbf{Bode's integral} (internally stable, $sL\to 0$)
\begin{align*}
\int_0^\infty \log|S(\ii\ww)|\,\dee\ww = \pi\sum p_k
\end{align*}
$p_k$: RHP poles of $L$. Pushing $|S|$ down somewhere pushes it up elsewhere: the \alert{waterbed effect}.
\par\medskip
\textbf{Nonlinear limits}
\begin{itemize}
\item Actuator amplitude/rate limits (minimum-time moves, windup)
\item Measurement noise (actuator wear, saturation)
\item Friction (limit cycles in balance systems)
\end{itemize}
\end{columns}
\LINKS{Reachability (T4), RHP zeros and delays (T5, T6), windup (T7), sensitivity (T8).}{Practice: the pracs and the project rig hit most of these limits.}
\end{frame}

%%%%%%%%%%%%%%%%%%%%%%%%%%%%%%%%
%% SYNTHESIS
%%%%%%%%%%%%%%%%%%%%%%%%%%%%%%%%

\gdef\currMODULE{Revision}
\gdef\currTOPIC{Putting it all together}
\section{Putting it all together}

\begin{frame}[plain]
\centering
\vfill
{\usebeamerfont{part title}\usebeamercolor[fg]{part title}Putting it all together\par}
\bigskip
\COURSEMAP{1,2,3,4,5,6,7,8,9}
\vfill
\end{frame}

\begin{frame}{Same idea, three languages}
\begin{center}
\scriptsize
\renewcommand\arraystretch{1.5}
\setlength\tabcolsep{4pt}\begin{tabular}{@{}>{\raggedright}p{1.6cm}>{\raggedright}p{2.8cm}>{\raggedright}p{2.6cm}>{\raggedright\arraybackslash}p{2.9cm}@{}}
\toprule
\textbf{Idea} & \textbf{State space} & \textbf{Transfer function} & \textbf{Frequency response}\\
\midrule
Stability & $\Real\lambda_j(A)<0$ & poles in LHP & Nyquist: $n_Z = 0$\\
Steady state & $D-CA^{-1}B$ & $G(0)$, final value thm & low-frequency gain\\
Speed & dominant eigenvalue, $T_s\approx 4/(\zz\ww_0)$ & dominant pole & bandwidth, $T_r\wb\approx 3$\\
Damping & $\zz = \cos\varphi$ (eigenvalue angle) & $\zz$ of dominant poles & resonant peak, phase margin\\
Zero SS error & augment $z=\int(y-r)$ & pole at $s=0$ in $C$ & $|L|\to\infty$ as $\ww\to 0$\\
Noise immunity & slower observer poles & roll-off in $C$ & small $|T|$, $|CS|$ at high $\ww$\\
Robustness & robust pole placement rules & avoid RHP cancellation & $\gainmargin$, $\phasemargin$, $\Msens$, $|\delta|<1/|T|$\\
\bottomrule
\end{tabular}
\end{center}
\medskip
\footnotesize
Being able to translate between these columns is what links the course together, and is what exam questions spanning several topics test.
\end{frame}

\begin{frame}{Running examples through the course}
\begin{center}
\scriptsize
\renewcommand\arraystretch{1.45}
\begin{tabular}{@{}p{2.1cm}p{9.0cm}@{}}
\toprule
\textbf{Example} & \textbf{Where it appears}\\
\midrule
Mass--spring--damper & T1 state model and second-order relations $\to$ T2 eigenvalues $\to$ T3 matrix exponential and coordinate changes $\to$ T5 $\frac{1}{ms^2+cs+k}$, DC gain $1/k$ by final value theorem\\
Cruise control & T2 nonlinear model with PI controller $\to$ T4 integral action ($\kkp=\frac{a_1-a}{b}$, $\kki=\frac{a_2}{b}$) $\to$ T7 windup on a steep hill $\to$ T9 uncertain vehicle mass\\
Balance system / inverted pendulum & T1 Segway model $\to$ T2 equilibria $\to$ T3 linearisation $\to$ T4 reachability and state feedback $\to$ T9 friction-induced oscillation\\
Vehicle steering & T4 state feedback design $\to$ T6 zero location and minimum phase $\to$ T9 RHP zero when reversing, slow zero and robust pole placement\\
Electric motor with delay & T6 critical delay $\tau_c$ $\to$ delay margin $\to$ T9 delays as a fundamental limit\\
\bottomrule
\end{tabular}
\end{center}
\medskip
\footnotesize
Revising one example end to end is a good way to rehearse the whole course.
\end{frame}

\begin{frame}{A design workflow}
\centering
\begin{tikzpicture}[node distance=3mm, >=Latex, thick, font=\scriptsize,
    wfstep/.style={draw=NorthTce, rounded corners=3pt, fill=NorthTce!8, text width=3.6cm, align=center, minimum height=7mm},
    note/.style={text width=5.2cm, align=left, font=\scriptsize}]
  \node[wfstep] (s1) {Model the plant};
  \node[wfstep, below=of s1] (s2) {Linearise at operating point};
  \node[wfstep, below=of s2] (s3) {Check structure};
  \node[wfstep, below=of s3] (s4) {Design controller};
  \node[wfstep, below=of s4] (s5) {Analyse the loop};
  \node[wfstep, below=of s5] (s6) {Check robustness \& limits};
  \node[wfstep, below=of s6] (s7) {Implement};
  \foreach \a/\b in {s1/s2,s2/s3,s3/s4,s4/s5,s5/s6,s6/s7} \draw[->] (\a) -- (\b);
  \draw[->, gray] (s6.west) -- ++(-0.5,0) |- node[pos=0.25, left, font=\tiny, text=gray]{iterate} (s4.west);
  \node[note, right=4mm of s1] {first principles or step/FRF data (T1, T5)};
  \node[note, right=4mm of s2] {Jacobian $A,B,C,D$; equilibria (T2, T3)};
  \node[note, right=4mm of s3] {stable? reachable? observable? RHP zeros, delays? (T2, T4, T5)};
  \node[note, right=4mm of s4] {PID tuning, pole placement + observer, or loop shaping (T4, T7, T8)};
  \node[note, right=4mm of s5] {Nyquist/Bode margins, Gang of Four, step response (T6, T8)};
  \node[note, right=4mm of s6] {$\Msens$, $|\delta|<1/|T|$, Bode integral, saturation (T9)};
  \node[note, right=4mm of s7] {filter D, anti-windup, setpoint weights, discretise (T1, T7)};
\end{tikzpicture}
\end{frame}

\begin{frame}{Which tool answers which question?}
\begin{center}
\scriptsize
\renewcommand\arraystretch{1.35}
\begin{tabular}{@{}p{4.6cm}p{4.9cm}l@{}}
\toprule
\textbf{Question} & \textbf{Tool} & \textbf{Topic}\\
\midrule
Is this equilibrium stable? & eigenvalues of the Jacobian $A$ & T2, T3\\
What is the response to an input or initial condition? & convolution equation, $\ee^{At}$ & T3\\
Can I steer the state anywhere? & rank of $W_r$ & T4\\
Can I reconstruct the state from $y$? & rank of $\Wob$ & T4\\
Where does a step response settle? & $G(0)$, final value theorem & T5\\
Is the closed loop stable, from open-loop data? & Nyquist: $n_Z = n_P + n_N$ & T6\\
How close is it to instability? & $\gainmargin$, $\phasemargin$, $\stabilitymargin = 1/\Msens$ & T6, T8\\
Quick controller with little model knowledge? & PID with ZN/CHR/Cohen--Coon/relay tuning & T7\\
Trade-off between disturbance and noise? & $S$, $T$, Gang of Four, loop shape & T8\\
What does one gain do to the poles? & root locus & T8\\
Will it still work if the model is wrong? & $|\delta|<1/|T|$, $\frac{\dee G}{G} = S\frac{\dee P}{P}$ & T9\\
Is the performance fundamentally limited? & RHP poles/zeros, delays, Bode integral, saturation & T9\\
\bottomrule
\end{tabular}
\end{center}
\end{frame}

\begin{frame}{Formula sheet (1/2): Modules I and II}
\scriptsize
\begin{columns}[t]
\column{0.5\linewidth}
\textbf{Models and response}
\begin{align*}
&\dot x = Ax+Bu,\quad y=Cx+Du\\
&x(t) = \ee^{At}x(0) + \textstyle\int_0^t \ee^{A(t-\tau)}Bu(\tau)\,\dee\tau\\
&A = \PDeriv{f}{x}\Big|_e,\; B = \PDeriv{f}{u}\Big|_e,\; C = \PDeriv{h}{x}\Big|_e,\; D = \PDeriv{h}{u}\Big|_e\\
&\wn = \sqrt{k/m},\quad \zz = c/(2\sqrt{km}),\quad \wdamp = \wn\sqrt{1-\zz^2}\\
&M_p = \ee^{-\pi\zz/\sqrt{1-\zz^2}},\quad T_s\approx 4/(\zz\wn)
\end{align*}
\textbf{State feedback and observers}
\begin{align*}
&W_r = [B\;\, AB\;\cdots\; A^{n-1}B],\quad \Wob = [C;\, CA;\, \cdots;\, CA^{n-1}]\\
&u = -Kx + k_f r,\quad \dot x = (A-BK)x + Bk_f r\\
&k_f = -1/\bigl(C(A-BK)^{-1}B\bigr)\\
&\dot{\hat x} = A\hat x + Bu + L(y - C\hat x),\quad \dot{\tilde x} = (A-LC)\tilde x
\end{align*}
\column{0.5\linewidth}
\textbf{Transfer functions}
\begin{align*}
&G(s) = C(sI-A)^{-1}B + D = b(s)/a(s)\\
&G(0) = D - CA^{-1}B = b_m/a_n\\
&\lim_{t\to\infty} f(t) = \lim_{s\to 0} sF(s),\quad \lim_{t\to 0} f(t) = \lim_{s\to\infty} sF(s)\\
&\frac{1}{ms^2+cs+k},\qquad \frac{K\wn^2}{s^2+2\zz\wn s+\wn^2}\\
&\ee^{-\tau s}\approx \frac{1-\tau s/2}{1+\tau s/2},\quad \arg \ee^{-\ii\ww\tau} = -\ww\tau
\end{align*}
\textbf{Frequency domain analysis}
\begin{align*}
&L = PC,\quad n_Z = n_P + n_N\\
&\gainmargin = 1/|L(\ii\ww_{\text{pc}})|,\quad \phasemargin = \ang{180} + \arg L(\ii\ww\gaincrossover)\\
&\phasemargin = \ww\gaincrossover\delaymargin,\quad \stabilitymargin = \min_\ww|1+L(\ii\ww)|
\end{align*}
\end{columns}
\end{frame}

\begin{frame}{Formula sheet (2/2): Module III}
\scriptsize
\begin{columns}[t]
\column{0.5\linewidth}
\textbf{PID}
\begin{align*}
&C(s) = \kkp + \frac{\kki}{s} + \frac{\kkd s}{1+s\TTf},\quad \TTi = \frac{\kkp}{\kki},\; \TTd = \frac{\kkd}{\kkp}\\
&\TTf = \TTd/N,\quad 5\le N\le 20\\
&P=\tfrac{b}{s+a},\;\text{PI}:\; \kkp = \tfrac{a_1-a}{b},\; \kki = \tfrac{a_2}{b}\\
&\text{ITD: } \tfrac{a}{\tau s}\ee^{-\tau s},\quad \text{FOTD: } \tfrac{K}{1+sT}\ee^{-\tau s},\quad \tau_n = \tfrac{\tau}{\tau+T}\\
&\text{relay: } k_c = \tfrac{4d}{\pi a},\quad \ww_c = 2\pi/T_c\\
&I(t_{k+1}) = I(t_k) + \kki h e(t_k) + \tfrac{h}{T_{\text{aw}}}(\mathrm{sat}(u)-u)
\end{align*}
\textbf{Root locus}
\begin{align*}
&a(s) + k\,b(s) = 0,\quad \text{asymptote angles } \tfrac{\pi + 2l\pi}{\npoleexcess}
\end{align*}
\column{0.5\linewidth}
\textbf{Sensitivity and loop shaping}
\begin{align*}
&S = \tfrac{1}{1+PC},\; T = \tfrac{PC}{1+PC},\; PS,\; CS,\quad S + T = 1\\
&\Msens = \max_\ww|S(\ii\ww)| = 1/\stabilitymargin\\
&T_r\wb\approx 3,\quad n\gaincrossover\approx -2 + \phasemargin/\ang{90}\\
&C = k\tfrac{s+a}{s+b}:\; \text{lead } a<b,\; \text{lag } a>b
\end{align*}
\textbf{Robustness and limits}
\begin{align*}
&\delta = \Delta/P,\quad |\delta(\ii\ww)| < 1/|T(\ii\ww)|\\
&\frac{\dee G_{yr}}{G_{yr}} = S\frac{\dee P}{P},\quad \frac{\dee G_{uw}}{G_{uw}} = -T\frac{\dee P}{P}\\
&\operatorname{rank}[A-\lambda I\;\; B] = n\quad\text{(PBH)}\\
&\textstyle\int_0^\infty \log|S(\ii\ww)|\,\dee\ww = \pi\sum p_k
\end{align*}
\end{columns}
\end{frame}

\begin{frame}{Common pitfalls}
\footnotesize
\begin{columns}[t]
\column{0.5\linewidth}
\begin{itemize}
\item \textbf{Units}: rad/s vs Hz; degrees vs radians (delay margin needs $\phasemargin$ in rad)
\item \textbf{dB}: gain margin as a ratio vs $20\log_{10}$
\item \textbf{Nyquist}: count encirclements of $-1$, not the origin; $n_P$ counts \emph{RHP} poles of $L$; clockwise is positive
\item \textbf{Cancellation}: never cancel a RHP pole or zero; check internal stability (all of the Gang of Four)
\item \textbf{Linearise about the equilibrium} (shift coordinates first)
\item \textbf{Transforms}: notes use both $x=Tz$ (T3 Jordan form) and $z=Tx$ (T3 coordinate change, T4 canonical forms): check which
\end{itemize}
\column{0.5\linewidth}
\begin{itemize}
\item \textbf{Frequencies}: resonance $\wres = \wn\sqrt{1-2\zz^2}$ $\neq$ damped $\wdamp = \wn\sqrt{1-\zz^2}$
\item \textbf{Rank}: $W_r$, $\Wob$ need rank $n$, not just `nonzero'
\item \textbf{Integral action} needs $b_n\neq 0$ (no zero at $s=0$) for the augmented system to be reachable
\item \textbf{Pure derivative} is improper: always filter
\item \textbf{Saturation}: linear analysis ignores it; check for windup
\item \textbf{Open-loop similarity} doesn't imply closed-loop similarity (T9)
\end{itemize}
\end{columns}
\end{frame}

%%%%%%%%%%%%%%%%%%%%%%%%%%%%%%%%
%% SELF TEST
%%%%%%%%%%%%%%%%%%%%%%%%%%%%%%%%

\begin{frame}{Self-test (1/3): Module I}
\footnotesize
\begin{enumerate}
\item For $\ddot x + 6\dot x + 5x = f$, find the poles, $\zz$, and the steady-state response to a unit step force.
  \uncover<2->{\par\alert{$s^2+6s+5=(s+1)(s+5)$: poles $-1$, $-5$. $\zz = 6/(2\sqrt5)\approx 1.34$ (overdamped). $x(\infty) = G(0) = 1/5$.}}
\item The pendulum $\dot x_1 = x_2$, $\dot x_2 = \sin x_1 - cx_2$ has equilibria at $x_1 = 0$ and $x_1 = \pi$. Which is stable?
  \uncover<3->{\par\alert{At $0$: $A=\Matr{0&1\\1&-c}$, $\lambda^2+c\lambda-1=0$ has a positive root: unstable (upright). At $\pi$: $A=\Matr{0&1\\-1&-c}$, stable for $c>0$.}}
\item Why must a model be shifted to its equilibrium before writing $\dot x = Ax+Bu$?
  \uncover<4->{\par\alert{A linear model has its only equilibrium at the origin; a constant offset makes it affine, not linear (T3).}}
\end{enumerate}
\end{frame}

\begin{frame}{Self-test (2/3): Module II}
\footnotesize
\begin{enumerate}
\item Is $A = \Matr{0&1\\0&0}$, $B=\Matr{0\\1}$ reachable? Design $K$ for closed-loop poles at $-2\pm 2\ii$.
  \uncover<2->{\par\alert{$W_r = \Matr{0&1\\1&0}$, rank 2: reachable. $\det(sI-A+BK) = s^2 + k_2 s + k_1 = s^2+4s+8$, so $K = \Matr{8 & 4}$.}}
\item For $L = k/\bigl(s(s+1)^2\bigr)$, for which $k$ is the closed loop stable?
  \uncover<3->{\par\alert{$\arg L = -\ang{180}$ at $\ww=1$ where $L(\ii)=-k/2$: stable for $0<k<2$, so $\gainmargin = 2/k$.}}
\item A loop has $\ww\gaincrossover = \SI{5}{rad/s}$ and $\phasemargin = \ang{45}$. Will an extra \SI{0.1}{s} delay destabilise it?
  \uncover<4->{\par\alert{$\delaymargin = \phasemargin/\ww\gaincrossover = (\pi/4)/5 \approx \SI{0.157}{s} > \SI{0.1}{s}$: still stable, with reduced margin.}}
\item Why does the observer not change the closed-loop poles set by $K$?
  \uncover<5->{\par\alert{The $(x,\tilde x)$ dynamics are block triangular: $p(s) = \det(sI-A+BK)\det(sI-A+LC)$.}}
\end{enumerate}
\end{frame}

\begin{frame}{Self-test (3/3): Module III}
\footnotesize
\begin{enumerate}
\item Design a PI controller for $P = b/(s+a)$ giving closed-loop poles with $\zz$ and $\ww_0$.
  \uncover<2->{\par\alert{$s^2 + (a+b\kkp)s + b\kki = s^2 + 2\zz\ww_0 s + \ww_0^2$: $\kkp = (2\zz\ww_0 - a)/b$, $\kki = \ww_0^2/b$.}}
\item A relay with $d=1$ produces a process oscillation of amplitude $0.5$. What is the critical gain?
  \uncover<3->{\par\alert{$k_c = 4d/(\pi a) = 4/(0.5\pi)\approx 2.55$.}}
\item $L = 3/(s+1)^3$ has $\stabilitymargin = 0.464$. What is the peak sensitivity, and what does it mean?
  \uncover<4->{\par\alert{$\Msens = 1/0.464\approx 2.2$: disturbances near $\ww_{\text{ms}}$ are amplified about $2.2\times$ by feedback.}}
\item With integral action, what is the steady-state output after a step load disturbance?
  \uncover<5->{\par\alert{$G_{yv}\approx s/\kki$ at low $\ww$, so by the final value theorem $y(\infty) = \lim_{s\to0} s\cdot\tfrac{s}{\kki}\cdot\tfrac{1}{s} = 0$.}}
\item $L$ has one RHP pole at $s=2$. How does that change the Bode integral, and why does it matter?
  \uncover<6->{\par\alert{$\int_0^\infty\log|S|\,\dee\ww = 2\pi > 0$: more sensitivity amplification overall, so a worse waterbed effect.}}
\end{enumerate}
\end{frame}

\begin{frame}{Closing}
\begin{itemize}
\item Module I: \emph{model} the system and understand its behaviour through equilibria and eigenvalues
\item Module II: \emph{analyse} feedback, in state space (pole placement, observers) and in frequency (transfer functions, Bode, Nyquist, margins)
\item Module III: \emph{design} with PID and loop shaping, and respect the limits set by sensitivity, uncertainty, RHP zeros, delays and saturation
\end{itemize}
\bigskip
\begin{center}
\emph{Stability and performance are linked:\\ every design is a compromise.}
\end{center}
\bigskip
\footnotesize
Next steps for revision: work through the worked examples for each concept, redo the workshop problems, and retrace one running example from model to implementation.
\end{frame}

\end{document}
